%%%%%%%%%%%%%%%%%%%%%%%%%%%%%%%%%%%%%%%%%%%%%%%%%%%%%%%%%%%%%%%%%%%%%%%%%
% Short Sectioned Assignment
% LaTeX Template
% Version 1.0 (5/5/12)
%
% This template has been downloaded from:
% http://www.LaTeXTemplates.com
%
% Original author:
% Frits Wenneker (http://www.howtotex.com)
%
% License:
% CC BY-NC-SA 3.0 (http://creativecommons.org/licenses/by-nc-sa/3.0/)
%
%%%%%%%%%%%%%%%%%%%%%%%%%%%%%%%%%%%%%%%%%%%%%%%%%%%%%%%%%%%%%%%%%%%%%%%%%%

%----------------------------------------------------------------------------------------
%	PACKAGES AND OTHER DOCUMENT CONFIGURATIONS
%----------------------------------------------------------------------------------------

\documentclass[paper=a4, fontsize=11pt]{scrartcl} % A4 paper and 11pt font size

\usepackage[T1]{fontenc} % Use 8-bit encoding that has 256 glyphs
\usepackage{fourier} % Use the Adobe Utopia font for the document - comment this line to return to the LaTeX default
\usepackage[english]{babel} % English language/hyphenation
\usepackage{amsmath,amsfonts,amsthm} % Math packages

\usepackage{lipsum} % Used for inserting dummy 'Lorem ipsum' text into the template

\usepackage{sectsty} % Allows customizing section commands
\allsectionsfont{\centering \normalfont\scshape} % Make all sections centered, the default font and small caps

\usepackage{fancyhdr} % Custom headers and footers
\pagestyle{fancyplain} % Makes all pages in the document conform to the custom headers and footers

%%%%my packages%%%%%%%%%%%%%%%%
\usepackage{enumerate}
\usepackage{graphicx}
\usepackage{fullpage}
\usepackage{xcolor}
\usepackage[linktocpage]{hyperref}
%%%%%%%%%%%%%%%%%%%%%%%%%%%%%%%%
\fancyhead{} % No page header - if you want one, create it in the same way as the footers below
\fancyfoot[L]{} % Empty left footer
\fancyfoot[C]{} % Empty center footer
\fancyfoot[R]{\thepage} % Page numbering for right footer
\renewcommand{\headrulewidth}{0pt} % Remove header underlines
\renewcommand{\footrulewidth}{0pt} % Remove footer underlines
\setlength{\headheight}{13.6pt} % Customize the height of the header

\numberwithin{equation}{section} % Number equations within sections (i.e. 1.1, 1.2, 2.1, 2.2 instead of 1, 2, 3, 4)
\numberwithin{figure}{section} % Number figures within sections (i.e. 1.1, 1.2, 2.1, 2.2 instead of 1, 2, 3, 4)
\numberwithin{table}{section} % Number tables within sections (i.e. 1.1, 1.2, 2.1, 2.2 instead of 1, 2, 3, 4)

\setlength\parindent{0pt} % Removes all indentation from paragraphs - comment this line for an assignment with lots of text

%----------------------------------------------------------------------------------------
%	TITLE SECTION
%----------------------------------------------------------------------------------------

\newcommand{\horrule}[1]{\rule{\linewidth}{#1}} % Create horizontal rule command with 1 argument of height

\title{	
\normalfont \normalsize 
%\includegraphics[scale=0.25]{images/AIMS_logo}
\textsc{African Institute for Mathematical Sciences -- \href{http://www.aims.ac.za}{AIMS}} 
\\ [25pt] % Your university, school and/or department name(s)
\includegraphics[scale=0.25]{images/logo_aims}
\horrule{0.5pt} \\[0.4cm] % Thin top horizontal rule
\huge Quantum Computing: Assignment 1 -- Linear Algebra \& Foundations of Quantum Mechanics \\ % The assignment title
\horrule{2pt} \\[0.5cm] % Thick bottom horizontal rule
}

\author{Evans Doe Ocansey \\ \textcolor{blue}{evans@aims.ac.za}} % Your name

\date{\normalsize\today} % Today's date or a custom date

\begin{document}

\maketitle % Print the title

%----------------------------------------------------------------------------------------
%	PROBLEM 1
%----------------------------------------------------------------------------------------

\section{Adjoint operators -- Sheet one, Question 4}
Consider a finite-dimensional vector space $V$ equipped with an inner product $(\cdot, \cdot)$. Suppose $A$ and $B$ are operators in $V$. Show the following identities
\begin{enumerate}[(a)]
\item $(A^{\dagger})^{\dagger} = A$. 
\item $(AB)^{\dagger} = B^{\dagger}A^{\dagger}$.
\item If $A$ and $B$ are both Hermitian operators then the commutator $C = [A, B]$ is anti-Hermitian i.e. $C^{\dagger} = -C$.
\end{enumerate}
%------------------------------------------------

\subsection{Solution -- Adjoint operators: Sheet one, Question 4}
Write your solution here
%----------------------------------------------------------------------------------------
%	PROBLEM 2
%----------------------------------------------------------------------------------------
\section{Projection operators -- Sheet one, Question 5}
%------------------------------------------------
Type  the question here

%------------------------------------------------

\subsection{Solution -- Projection operators: Sheet one, Question 5}

Write your solution here

%----------------------------------------------------------------------------------------
%	PROBLEM 3
%----------------------------------------------------------------------------------------

\section{Measurements and expectation values -- Sheet two, Question 1}
%------------------------------------------------
Type  the question here

%------------------------------------------------

\subsection{Solution -- Projection operators: Sheet two, Question 1}

Write your solution here

%----------------------------------------------------------------------------------------
%	PROBLEM 4
%----------------------------------------------------------------------------------------

\section{Subspaces of a vector space -- Sheet two, Question 3}
%------------------------------------------------
Type  the question here

%------------------------------------------------

\subsection{Solution -- Subspaces of a vector space: Sheet two, Question 3}

Write your solution here
\end{document}